\section{Задача и мотивация}

Открытословарные семантические карты (ConceptGraphs~\cite{conceptgraphs}, BBQ, OpenScene) позволяют роботу искать объекты и избегать их по произвольному текстовому запросу: «не подъезжай к растению», «держись дальше от велосипеда». Такая карта собирается из двух ненадёжных источников. Метки дают детектор или сегментатор открытого словаря, и они ошибаются. Положение объектов задаётся оценкой позы робота (SLAM, одометрия), и она дрейфует. Планировщик же обычно считает карту правдой: ищет путь по предсказанным областям классов и держит до них фиксированный зазор.

Цена такого допущения не измеряется. Бенчмарк лаборатории BE2R OSMa-Bench~\cite{osmabench} показывает, что качество открытословарных карт сильно деградирует: при динамическом освещении f-mIoU ConceptGraphs падает с 28{,}00 до 14{,}29 (табл.~I в~\cite{osmabench}). Но бенчмарк оценивает карту как изображение, без задачи движения. Обзор 3D-графов сцены 2026~г.~\cite{survey3dsg} формулирует то же в общем виде: влияние ошибок восприятия на поведение робота почти никогда не измеряют напрямую.

Есть и работы, которые дают планировщику статистическую гарантию поверх ненадёжного восприятия, с помощью конформного предсказания (conformal prediction, CP). Sundarsingh и др.~\cite{sundarsingh} калибруют множества меток клеток семантической карты, а Perceive With Confidence~\cite{pwc} калибрует раздувание карты занятости. Обе работы явно предполагают, что \emph{поза робота известна}. Работа о композиции гарантий~\cite{baral} показывает, что раздельно откалиброванные модули не дают откалиброванной системы.

\textbf{Задача работы} --- измерить, как совместные ошибки метки и позы проходят в нарушения безопасности при навигации по открытословарной карте, и проверить, может ли \emph{одна} калиброванная величина --- «расстояние промаха» --- дать планировщику гарантию сразу на оба источника ошибки.

Работа выполнена как тестовое задание лаборатории BE2R (направление «семантическое картирование, визуальная привязка и навигация») и одновременно как первый шаг ВКР по планированию движения в условиях неопределённости (научный руководитель И.~С.~Довгополик). Карта и её ошибки здесь --- источник неопределённости, а предмет работы --- решение планировщика: куда можно ехать и с каким запасом.
